\documentclass[10pt,a4paper]{amsart}
\usepackage[margin=19mm]{geometry}
\usepackage{amsmath,amssymb,mathtools}
\usepackage[scr=rsfs,cal=euler]{mathalfa}
\usepackage{xfrac}
\usepackage[unicode,hidelinks,bookmarks=false]{hyperref}
\newcommand{\ie}{i.e.\ }
\newcommand{\cf}{cf.\ }
\newcommand{\resp}{respectively\ }
\newcommand{\Subsection}{\subsection}
\providecommand{\nobreakdash}{\nobreak\hbox{-}}
\setlength{\emergencystretch}{2em}
\multlinegap0pt
\usepackage[all]{xy}
\usepackage{amssymb}
\usepackage{bm}
\usepackage[shortlabels]{enumitem}
\newtheorem{lem}{Lemma}
\newtheorem{prop}{Proposition}
\newcommand{\A}{{\mathbb A}}
\newcommand{\pone}{{\mathbb P}^1}
\title{{\bf Exotic Hopf maps, weight shifting and applications to vector bundles}}
\author{Jean Fasel, William Hornslien}
\date{Source: arXiv:2604.12541v1 (CC BY 4.0)}
\begin{document}
\maketitle

\noindent\emph{Source and licence.} {\bf Exotic Hopf maps, weight shifting and applications to vector bundles}. Version arXiv:2604.12541v1 (CC BY 4.0), distributed by arXiv under the Creative Commons Attribution 4.0 International licence, http://creativecommons.org/licenses/by/4.0/, as recorded in the arXiv licence metadata for this exact version and shown on the abstract page https://arxiv.org/abs/2604.12541v1. Full licence notice: \texttt{licenses/arxiv-2604.12541-CC-BY-4.0.txt}.\par\smallskip

\noindent\textit{Editorial scope.} Counted unit: the article's prop prop:etasusp (source0.tex lines 772--777). The article's complete original proof of it is delivered with it (source0.tex lines 778--786, one \texttt{proof} environment of 9 lines). No mathematics is written, repaired or re-derived: the statement and the proof are the article's own lines, in the article's own notation, and no retained line is substituted or re-flowed.  Retained uncounted as the source-local context the proof needs: lines 670--673.  Nothing was omitted from any retained span, so the package carries no omission record.  No retained line contains a citation, so the wrapper carries no bibliography and no bibliography entry.  No retained line contains a figure environment, an image-inclusion command or drawing code, so no image is delivered and the package declares no asset dependency.  Editorial formatting only, in addition: an AMS \texttt{amsart} preamble that restates the article's own declarations and macros used by the retained lines, namely the environments \texttt{lem}, \texttt{prop} on separate counters, together with the standard packages those lines need.  The retained environments print in this document as Lemma 1, Proposition 1 in order of appearance, whereas the article prints the same items with the numbering of its own full document.  The complete native source of the article as distributed in the arXiv e-print for this exact version is bundled unchanged as \texttt{sources/198408/source0.tex}.
\begin{lem}
\label{lem:ponesuspoddeven}
The $\pone$-suspension of a map $f:Q_{2n-1}\to Q_{2m}$ (resp.$f:Q_{2n}\to Q_{2m})$  given by the ideal $\langle f_1, \ldots, f_n, f_z \rangle $ is the map $F:Q_{2n+1}\to Q_{2m+2}$ (resp. $F:Q_{2n+2}\to Q_{2m+2}$) given by the ideal $\langle f_1, \ldots, f_n, x_{n+1},f_z \rangle$.
\end{lem}

\begin{prop}
\label{prop:etasusp}
Let $n\geq 1$. The map $f\colon Q_{2n+1} \to Q_{2n}$ defined by the ideal $(x_3, \ldots, x_n, x_1x_2)$ is the $(n-1)$-fold $\mathbb{P}^1$-suspension of the Hopf map $\eta \colon Q_3 \to Q_2.$


\end{prop}

\begin{proof}
We first need to prove that the claimed map is the Hopf map when $n=1$. Recall that we have a homotopy equivalence $\pi \colon Q_3 \to \A^2 \smallsetminus\{0\}$ sending $(x_1,x_2,y_1,y_2)$ to the point $(x_1,x_2)$ in $\A^2 \smallsetminus\{0\}$. There is also a homotopy equivalence $Q_2 \to \pone$ which maps a point $(x_1,y_1,z) \in Q_2$ to the point $[z:x_1]$ or $[y_1 : 1-z]$ (whichever is nonzero). The Hopf map $\eta\colon \A^2 \smallsetminus\{0\} \to \pone$ is defined by $(x_1,x_2) \to [x_1:x_2]$. It is straightforward to check that the square 
\begin{equation*}
    \xymatrix{Q_3\ar[r]^\eta\ar[d]^\pi & Q_2\ar[d]^\pi \\
\A^2 \smallsetminus\{0\} \ar[r]_{\eta} & \mathbb{P}^1.}
\end{equation*}
commutes. The claim then follows from Lemma \ref{lem:ponesuspoddeven}.

\end{proof}

\end{document}
